% Context from chapters/shannon.tex; for mathematical reference, not compilation.
e Sobolev regularity to control approximation and denoising errors.

   
\paragraph{Fourier series.}

Let $\TT=\RR/(2\pi\ZZ)$ denote the circle of length $2\pi$.
 
A function $f\in L^2(\TT)$ is identified with a $2\pi$-periodic function on $\RR$. Its Fourier coefficients are
\eq{
	\foralls k \in \ZZ, \quad 
	\hat f_k \eqdef \frac{1}{2\pi}\int_0^{2\pi} f(x) e^{-\imath x k} \d x.
}
Equivalently, $\hat f_k = \dotp{f}{e_k}$, where $\dotp{f}{g} \eqdef \frac{1}{2\pi} \int_\TT f(x) \bar g(x) \d x$ and $e_k(x) \eqdef e^{\imath x k}$.
 
The functions $(e_k)_{k\in\ZZ}$ form an orthonormal basis for this inner product, giving the expansion
\eql{\label{eq-fourier-series}
	f = \sum_{k \in \ZZ} \dotp{f}{e_k} e_k
}
which means $\norm{f-\sum_{k=-N}^N \dotp{f}{e_k} e_k}_{L^2(\TT)} \rightarrow 0$ for $N \rightarrow +\infty$.
 
For a differentiable function, the series~\eqref{eq-fourier-series} converges pointwise at every $x \in \TT$. If $f$ is of class $C^2$ on $\TT$, then $\hat f_k = O(1/k^2)$, which gives normal convergence and hence uniform convergence.
 
For a piecewise $C^1$ function, the Fourier series converges at a jump to the average of the one-sided limits. Gibbs oscillations prevent uniform convergence across the jump.


   
\paragraph{Poisson formula.}

The Poisson formula relates sampling in the signal domain to periodization in the Fourier domain.
 
For a function $h(\om)$ on $\RR$, define its periodization by
\eql{\label{eq-periodizing}
	h_P(\om) \eqdef \sum_n h(\om-2\pi n).
} 
If $h \in L^1(\RR)$, the series converges absolutely almost everywhere and $\norm{h_P}_{L^1(\TT)} \leq \norm{h}_{L^1(\RR)}$, with equality for nonnegative functions. Here the norm on $\TT$ uses unnormalized Lebesgue measure over one period. Indeed,
$$
	|h_P(x)| \leq (|h|_P)(x), \quad\Longrightarrow\quad
	\norm{h_P}_{L^1} \leq  \norm{|h|_P}_{L^1} = \norm{h}_{L^1}. 
$$
 
For $h=\hat f$, Proposition~\ref{prop-poisson} states that the following diagram
\eq{
	\begin{array}{rcccl}
						& f(x)  &  \overset{\Ff}{\longrightarrow} &  \hat f(\om) &\\
		\text{sampling}& \downarrow & & \downarrow &\text{periodization} \\
						& (f(n))_n  &  \overset{\text{Fourier series}}{\longrightarrow} &  \sum_n f(n) e^{-\imath \om n} &\\
	\end{array}
}
is commutative. The minus sign in $\sum_n f(n)e^{-\imath\om n}$ is opposite to the sign in Fourier series synthesis; this accounts for the index reversal in the proof below.

\begin{prop}[Poisson formula]\label{prop-poisson}
Assume that $f\in C(\RR)$, that $\hat f$ has compact support, and that $|f(x)| \leq C(1+|x|)^{-3}$ for some $C$. Then
\eql{\label{eq-poisson-formula}
	\foralls \om \in \RR, \quad
	\sum_n f(n) e^{-\imath \om n} = \hat f_P(\om).
}
\end{prop}
\begin{proof}
	The decay of $f$ makes $\hat f$ continuously differentiable. Its compact support makes the periodization locally a finite sum, so $(\hat f)_P$ is also $C^1$. It therefore equals its Fourier series:
	\eql{\label{eq-poisson-series}
		(\hat f)_P(\om) = \sum_k c_k e^{\imath k \om},
	} 
	where
	\begin{align*}
		c_k = \frac{1}{2\pi} \int_0^{2\pi} (\hat f)_P(\om) e^{-\imath k \om} \d \om = 
		\frac{1}{2\pi} \int_0^{2\pi} \sum_n \hat f(\om-2\pi n) e^{-\imath k \om}  \d \om .
	\end{align*}
	Absolute integrability follows from
	\eq{
		\int_0^{2\pi} \sum_n |\hat f(\om-2\pi n) e^{-\imath k \om}|  \d \om = \int_\RR |\hat f| 
	}
	which is finite because $\hat f$ is continuous and compactly supported. We may therefore exchange the sum and integral:
	\eq{
		c_k = \sum_n \frac{1}{2\pi} \int_0^{2\pi} \hat f(\om-2\pi  n) e^{-\imath k \om}  \d \om
		= \frac{1}{2\pi} \int_{\RR} \hat f(\om) e^{-\imath k \om}  \d \om
		= f(-k)
	}
	where we used the inverse Fourier transform formula~\eqref{eq-i-ft}, which is valid because $\hat f \in L^1(\RR)$.
\end{proof}

   
\paragraph{Shannon's sampling theorem.}

Shannon's sampling theorem gives a sufficient condition for reconstructing a signal from the samples $f \mapsto (f(ns))_n$ at spacing $s>0$.
 
The bandlimiting condition is $\supp(\hat f) \subset [-\pi/s,\pi/s]$. By Fourier inversion~\eqref{eq-i-ft}, such a signal is $C^\infty$ and has an entire extension to the complex plane.
 
Shannon's 1949 paper placed the sampling theorem at the heart of communication theory, while acknowledging earlier interpolation results and Nyquist's work on transmission rates.
 
Figure~\ref{fig-sampling-aliasing} illustrates the reconstruction argument and the aliasing that occurs when spectral replicas overlap; see also Figure~\ref{fig-aliasing}.

\begin{figure}[tbp]
\centering
\BookDrawing{tikz/shannon-sampling-aliasing}
\caption{Sampling and Shannon interpolation in the frequency and time domains, with unit sample spacing. Left: the bandlimiting condition of Theorem~\ref{thm-shannon-sampling} holds. Right: overlapping spectral replicas cause aliasing; the reconstructed signal still passes through every sample. The rows show the original signal, sampling, the interpolation kernel, and reconstruction.}
\label{fig-sampling-aliasing}
\end{figure}





\begin{thm} \label{thm-shannon-sampling}
	If $f\in C(\RR)$, $|f(x)| \leq C(1+|x|)^{-3}$ for some $C$ and $\supp(\hat f) \subset [-\pi/s,\pi/s]$, then one has
	\eql{\label{eq-shannong-interp}
		\foralls x \in \RR, \quad 
		f(x) = \sum_n f(n s) \sinc(x/s-n) \qwhereq
		\sinc(u) = \frac{\sin(\pi u)}{\pi u},\quad \sinc(0)=1
	}
	with uniform convergence.
\end{thm}

\begin{proof} 
	Set $g \eqdef f(s \cdot)$. Then $\hat g=1/s \hat f(\cdot/s)$, as follows from the substitution $z=s x$:
	\eq{
		\hat g(\om) = \int f(s x) e^{-\imath \om x} \d x = \frac{1}{s} \int f(z) e^{-\imath (\om/s) z} \d z = \hat f(\om/s)/s, 
	} 
	It is therefore enough to prove the result for unit sample spacing:
	$$
		g(x) = \sum_n g(n) \sinc(x-n), 
	$$
	For the rest of the proof, write $f$ in place of $g$.
	 
	The compact support hypothesis implies $\hat f(\om) = 1_{[-\pi,\pi]}(\om) \hat f_P(\om)$.  
	Combining the inversion formula~\eqref{eq-i-ft} with the Poisson formula~\eqref{eq-poisson-formula} gives
	\eq{
		f(x) = \frac{1}{2\pi} \int_{-\pi}^\pi \hat f_P(\om) e^{\imath \om x} \d \om
		= \frac{1}{2\pi} \int_{-\pi}^\pi \sum_n f(n) e^{\imath \om (x-n)} \d \om.
	} 
	The decay of $f$ gives $\int_{-\pi}^\pi \sum_n |f(n) e^{\imath \om (x-n)}| \d \om = 2\pi\sum_n |f(n)| < +\infty$. We may therefore exchange summation and integration:
	\eq{
		f(x) = \sum_n f(n)  \frac{1}{2\pi} \int_{-\pi}^\pi e^{\imath \om (x-n)} \d \om = \sum_n f(n) \sinc(x-n).
	}
The series converges uniformly by the Weierstrass test, since $|\sinc(x-n)|\leq 1$ and $\sum_n|f(n)|<\infty$.
\end{proof}

\begin{figure}[tbp]
\centering
\BookDrawing{tikz/shannon-sinc}
\caption{The sinc interpolation kernel.}
\label{fig-review-shannon-sinc}
\end{figure}
The sinc kernel decays slowly and oscillates, making it inconvenient for local interpolation. Oversampling gives more flexibility: if $\supp(\hat f)\subset[-\pi/s',\pi/s']$ with $s'>s$, the spectrum occupies a strict subinterval of the Nyquist interval. One can then choose a reconstruction 